\chapter{The pilot session, in full}
\label{app:pilot}

\noindent\emph{\Cref{sec:pilot-results} reports the headline pilot numbers.
This appendix is the full record: the coverage actually achieved, every
session, and all nine figures produced from the raw 20 Hz position traces and
clutch event logs.}

Sixteen sessions were recorded on 2026-09-03 and 2026-09-04, drawn from
\repo{recordings/sessions/}. Duration, hand-travel distance and lag are
computed from the full-resolution end-effector and controller position trace
(20 Hz); clutch and emergency-stop counts are computed from each session's
event log; participant identity was replaced with P1--P4 in every table and
figure below and nowhere in this appendix or its source data does a real name
appear, consistent with \repo{write\_manifest()} refusing to record one
(\cref{sec:ethics}).

\section{Coverage}

\begin{table}[htbp]
  \centering
  \small
  \caption[Pilot coverage]{Which participant performed which task under which
  condition. Coverage is uneven because a pilot session is limited by time; the
  pooled comparison of \cref{sec:pilot-results} is read against this table
  rather than against an assumed complete grid.}
  \label{tab:pilot-coverage}
  \begin{tabular}{@{}lccc@{}}
    \toprule
    Participant & Object tracking & Target reaching & Position matching \\
    \midrule
    P1 & both (direct $\times2$) & --- & --- \\
    P2 & both & both (shared $\times2$) & --- \\
    P3 & direct only & shared only & both \\
    P4 & both & --- & both \\
    \bottomrule
  \end{tabular}
\end{table}

\section{Every session}

\begin{table}[htbp]
  \centering
  \small
  \caption[Every pilot session]{Every recorded session. Duration includes
  setup and hesitation. Active share is the fraction of the session spent
  with the controller grip engaged. P1's two direct-control object-tracking
  sessions are a short warm-up followed by the recorded attempt, both
  retained as separate points.}
  \label{tab:pilot-sessions}
  \begin{tabular}{@{}llrrrrr@{}}
    \toprule
    P & Task & Cond. & Duration (s) & Hand travel (m) & Re-grips & Active share \\
    \midrule
    P1 & Object tracking & Direct & 147.4 & 0.31 & 19 & \SI{50}{\percent} \\
    P1 & Object tracking & Direct & 392.8 & 1.85 & 53 & \SI{53}{\percent} \\
    P1 & Object tracking & Shared & 114.9 & 1.50 & 12 & \SI{61}{\percent} \\
    P2 & Object tracking & Direct & 38.7  & 0.90 & 2  & \SI{61}{\percent} \\
    P2 & Object tracking & Shared & 169.6 & 1.57 & 25 & \SI{76}{\percent} \\
    P2 & Target reaching & Direct & 227.9 & 2.77 & 25 & \SI{52}{\percent} \\
    P2 & Target reaching & Shared & 78.6  & 0.70 & 11 & \SI{65}{\percent} \\
    P2 & Target reaching & Shared & 357.9 & 5.39 & 41 & \SI{73}{\percent} \\
    P3 & Object tracking & Direct & 617.4 & 7.09 & 52 & \SI{83}{\percent} \\
    P3 & Target reaching & Shared & 458.1 & 3.18 & 51 & \SI{74}{\percent} \\
    P3 & Position matching & Direct & 315.9 & 3.72 & 35 & \SI{67}{\percent} \\
    P3 & Position matching & Shared & 196.8 & 4.08 & 14 & \SI{85}{\percent} \\
    P4 & Object tracking & Direct & 195.6 & 2.07 & 37 & \SI{68}{\percent} \\
    P4 & Object tracking & Shared & 186.9 & 1.98 & 25 & \SI{42}{\percent} \\
    P4 & Position matching & Direct & 171.9 & 1.65 & 12 & \SI{30}{\percent} \\
    P4 & Position matching & Shared & 111.7 & 1.98 & 12 & \SI{40}{\percent} \\
    \bottomrule
  \end{tabular}
\end{table}

\section{Figures}

\begin{figure}[htbp]
  \centering
  \includegraphics[width=0.85\linewidth]{pilot/fig2_hand_distance.png}
  \caption[Pilot: hand-travel distance]{Distance the robot's hand travelled,
  by task and condition. Pooled distance is essentially unchanged between
  conditions, which is why \cref{sec:pilot-limits} reads it as one of the
  more reliable measures in a pilot where session duration is confounded with
  setup time.}
  \label{fig:pilot-distance}
\end{figure}

\begin{figure}[htbp]
  \centering
  \includegraphics[width=0.85\linewidth]{pilot/fig3_regrips.png}
  \caption[Pilot: re-grips per session]{How often the operator paused and
  re-engaged the controller grip, by task and condition.}
  \label{fig:pilot-regrips}
\end{figure}

\begin{figure}[htbp]
  \centering
  \includegraphics[width=0.85\linewidth]{pilot/fig4_active_share.png}
  \caption[Pilot: active-control share]{Share of the session spent with the
  grip engaged, by task and condition.}
  \label{fig:pilot-active}
\end{figure}

\begin{figure}[htbp]
  \centering
  \includegraphics[width=0.85\linewidth]{pilot/fig5_lag.png}
  \caption[Pilot: spatial lag]{Mean spatial lag the VR link itself reports
  between the controller pose and the commanded pose, by task and condition.
  Both conditions stay under \SI{5}{\milli\metre}.}
  \label{fig:pilot-lag}
\end{figure}

\begin{figure}[htbp]
  \centering
  \includegraphics[width=0.72\linewidth]{pilot/fig7_paired_regrips.png}
  \caption[Pilot: paired re-grips]{Re-grip count for the same six paired
  cases as \cref{fig:pilot-paired}.}
  \label{fig:pilot-paired-regrips}
\end{figure}

\begin{figure}[htbp]
  \centering
  \includegraphics[width=0.95\linewidth]{pilot/fig8_trajectory_comparison.png}
  \caption[Pilot: recorded end-effector path]{The actual recorded
  end-effector path in space for one participant on the same task under both
  conditions, recorded back to back. Neither path is smooth; a shorter,
  straighter line would indicate a more direct approach to the target.}
  \label{fig:pilot-trajectory}
\end{figure}

\begin{figure}[htbp]
  \centering
  \includegraphics[width=0.95\linewidth]{pilot/fig9_engagement_timeline.png}
  \caption[Pilot: engagement timeline]{When the operator was actively
  gripping the controller across the session, for the same participant and
  task as \cref{fig:pilot-trajectory}. Gaps are the operator releasing to
  reposition their hand before re-engaging.}
  \label{fig:pilot-timeline}
\end{figure}

\section{What this pilot does not establish}

Four participants with uneven task coverage is a pilot-scale sample: none of
the differences above should be read as a proven effect, and no significance
test is reported for that reason. The pooled figures mix three tasks
together; \cref{fig:pilot-paired,fig:pilot-paired-regrips} are the six cases
that isolate the condition within one participant and one task, and they are
the part of this pilot closest to the paired design of \cref{app:study}.
Session duration runs from the start of the recording to the end, including
any explanation or hesitation, so it is a session-length measure rather than
a pure motor-execution time; hand-travel distance and lag, drawn directly
from the position trace, are not subject to the same confound. Every
session's emergency stop was engaged once, near the end, to stop the
recording -- a procedural marker rather than a safety incident -- and is
excluded from every figure above. Most importantly, this pilot has no
separate wearer: every measure here is what the operator experienced, and
the wearer-side hypotheses of \cref{app:study} (H3, H4) remain untested.
